\documentclass[12pt, a4paper]{article}
\usepackage[utf8]{inputenc}
\usepackage[lithuanian]{babel}
\usepackage[T1]{fontenc}
\usepackage{amsmath, amssymb}
\usepackage{geometry}
\usepackage{enumitem}

\geometry{top=2cm, bottom=2cm, left=2cm, right=2cm}

\title{\textbf{Namų darbas nr. 1}}
\date{}
\begin{document}

\maketitle

\section{Sąlygos}

\begin{enumerate}[label=\textbf{\arabic* užduotis.}, leftmargin=*]
    % 1 lygis: Baziniai veiksmai su vienanariais
    \item Atlikite veiksmus su vienanariais ir nustatykite, ar gautas reiškinys yra vienanaris. Jei taip, nurodykite jo koeficientą ir laipsnį:
    \begin{enumerate}[label=\alph*)]
        \item $32a^6 b^{10} : (2a^2 b^3)^2$
        \item $\dfrac{5x^3 y z \cdot 4xyz^{10}}{8x^2 y^2 z^3}$
        \item $(3abc^4)^3 : (9a^2 b^2 c^{13})$
    \end{enumerate}

    % 2 lygis: Tapatybės / koeficientų sulyginimas (tiesinis lygčių radimas)
    \item Raskite nežinomų kintamųjų $a$, $b$ ir $c$ reikšmes, su kuriomis lygybė yra tapatybė:
    \[
    (2x + a)(bx + 4) = cx^2 + 14x + 8
    \]

    % 3 lygis: Tapatybės / koeficientų sulyginimas (su parametrais)
    \item Raskite tokias kintamųjų $p$, $q$ ir $r$ reikšmes, su kuriomis lygybė galioja su visomis $x$ reikšmėmis:
    \[
    (3x - p)(2x + q) = rx^2 + x - 10
    \]

    % 4 lygis: Sudėtinis algebrinis įrodymas taikant greitosios daugybos formules
    \item Įrodykite, kad su visomis leistinomis $x$ reikšmėmis ($x \neq 0$) trupmenos reikšmė yra pastovi:
    \[
    \frac{(4x - 1)(2x + 3) - 2(x - 2)^2 - (3x - 1)(3x + 1) + 3\left(x^2 + 3\frac{1}{3}\right)}{(3x + 1)(2x + 5) - 2(x + 2)^2 - (2x - 1)(2x + 1) + 3}
    \]
    Kodėl $x$ negali būti lygus nuliui?  
    % 5 lygis: Sudėtinis daugiapakopis skaitinis reiškinys su DBD, šaknimis ir laipsniais
    \item Apskaičiuokite reiškinio reikšmę:
    \[
    \frac{6\frac{1}{4} \cdot 2{,}8 - 3{,}4 : \frac{34}{55}}{(\sqrt{20} - 2\sqrt{45})^2} \cdot \frac{1}{\sqrt{\left(1\frac{9}{16}\right)^{-1}}} \cdot \sqrt[40]{4^{\text{DBD}(560, 880)}}
    \]
    \textit{Pastaba: $\text{DBD}(560, 880)$ – didžiausias bendrasis skaičių 560 ir 880 daliklis.}

    % 6 lygis: Abstrakčiausia algebrinė rodiklinė lygtis
    \item Išspręskite lygtį kintamojo $m$ atžvilgiu:
    \[
    \frac{(a^{m+1})^{m-3} \cdot a^{m^2 - 3m + 4}}{a^{3-2m}} = a^{6-3m}
    \]
\end{enumerate}

\newpage

\section{Atsakymai}

\begin{enumerate}[label=\textbf{\arabic* užduotis.}, leftmargin=*]
    \item 
    \begin{enumerate}[label=\alph*)]
        \item Yra vienanaris. Koeficientas: $8$, laipsnis: $6$.
        \item Yra vienanaris. Koeficientas: $\dfrac{5}{2}$ (arba $2{,}5$), laipsnis: $10$.
        \item Nėra vienanaris.
    \end{enumerate}

    \item $a = 2, \quad b = 3, \quad c = 6$

    \item $p = 5, \quad q = 2, \quad r = 6$

    \item Trupmenos reikšmė nepriklausomai nuo $x$ lygi $2$.

    \item $3$

    \item $m_1 = -2, \quad m_2 = 2$
\end{enumerate}

\end{document}